\chapter{State of the Art, October 2026}
\label{ch:sota}

This chapter surveys the six areas the project draws on. In each, the newest
verified work comes first. The search closed on 2~October~2026. Each table
gives the month a work first appeared, its venue where one exists, and the
result most relevant here. Older work appears only where it is still the
reference baseline. The aim is not an exhaustive review but a reliable map
of the frontier that the proposed method must be measured against.

\section{Real-time general detectors}
\label{sec:sota-rt}

The accuracy--latency frontier moved in two ways between 2024 and 2026. The
YOLO line kept improving its training recipe and heads. YOLOv10 removed the
need for NMS~\citep{wang2024yolov10}. YOLOv12 made attention efficient
enough for real time~\citep{tian2025yolov12}. YOLOv13 added hypergraph
correlation modules~\citep{lei2025yolov13}. YOLO26 removed distribution focal
loss and added a small-target-aware label assignment~\citep{yolo26_2026}.

At the same time, detection transformers began to draw on large
self-supervised vision models. RF-DETR uses a DINOv2 backbone with weight-sharing
architecture search~\citep{robinson2025rfdetr,oquab2024dinov2}. DEIMv2 and
RT-DETRv4 use DINOv3, as features or as a teacher during
training~\citep{deimv2_2025,rtdetrv4_2025,simeoni2025dinov3}. RF-DETR is the
first real-time detector reported above 60 COCO AP. Table~\ref{tab:sota-rt}
and Figure~\ref{fig:frontier} summarise the field.

\begin{table}[htbp]
\centering
\caption[Real-time detectors, newest first]{Real-time detectors, newest
first. COCO val2017 AP and latency on a T4 GPU with TensorRT FP16, as
reported by each work. The ranges run from the smallest to the largest
model.}
\label{tab:sota-rt}
\small
\begin{tabularx}{\textwidth}{@{}P{1.35cm}P{2.6cm}P{2.3cm}L@{}}
\hd{Date} & \hd{Model} & \hd{Venue} & \hd{COCO AP and latency} \\ \gap
06/2026 & YOLO26~\citep{yolo26_2026} & arXiv; released 01/2026 & 40.9 to 57.5 at 1.7 to 11.8\,ms; optional NMS-free head; small-target-aware assignment \\
11/2025 & RF-DETR~\citep{robinson2025rfdetr} & ICLR 2026 & 48.4 to 60.1 at 2.3 to 17.2\,ms; DINOv2 backbone \\
10/2025 & RT-DETRv4~\citep{rtdetrv4_2025} & ECCV 2026 & 49.8 to 57.0 at 3.7 to 12.9\,ms; DINOv3 teacher, no inference cost \\
09/2025 & DEIMv2~\citep{deimv2_2025} & arXiv & 38.5 (1.5M params) to 57.8; first model under 10M params above 50 AP \\
06/2025 & YOLOv13~\citep{lei2025yolov13} & arXiv & 41.6 to 54.8 at 2.0 to 14.7\,ms \\
02/2025 & YOLOv12~\citep{tian2025yolov12} & NeurIPS 2025 & 40.4 to 55.4 at 1.6 to 10.4\,ms; needs FlashAttention for full speed \\
12/2024 & DEIM~\citep{huang2025deim} & CVPR 2025 & 43.0 to 56.5; dense one-to-one matching, half the training time \\
10/2024 & D-FINE~\citep{peng2025dfine} & ICLR 2025 & 42.8 to 55.8; 59.3 with Objects365 pretraining \\
09/2024 & YOLO11~\citep{yolo11_2024} & software & 39.5 to 54.7 at 1.5 to 11.3\,ms \\
05/2024 & YOLOv10~\citep{wang2024yolov10} & NeurIPS 2024 & 38.5 to 54.4 at 1.8 to 10.7\,ms; first NMS-free YOLO \\
04/2023 & RT-DETR~\citep{zhao2024rtdetr} & CVPR 2024 & baseline of HPS-DETR; v2 and v3 followed in 2024 \\ \gap
\multicolumn{4}{@{}l}{\textit{Accuracy ceilings, not real-time}} \\
08/2025 & DINOv3 7B with Plain-DETR~\citep{simeoni2025dinov3} & arXiv & 65.6 AP\unv \\
11/2022 & Co-DETR, ViT-L~\citep{zong2023codetr} & ICCV 2023 & 66.0 AP on test-dev \\
\end{tabularx}
\end{table}

\begin{figure}[htbp]
\centering
\begin{tikzpicture}
  \begin{axis}[
      width=12.4cm, height=7.6cm,
      xmin=0, xmax=18.5, ymin=38, ymax=61.5,
      axis lines=left, axis line style={draw=cNeutral, line width=0.5pt},
      tick style={draw=cNeutral}, tick label style={font=\scriptsize, text=muted},
      xlabel={latency on T4, TensorRT FP16 (ms)}, ylabel={COCO val AP},
      label style={font=\footnotesize, text=muted},
      xtick={0,2,4,6,8,10,12,14,16,18}, ytick={40,45,50,55,60},
      legend style={draw=none, fill=none, font=\scriptsize, text=ink,
                    at={(0.98,0.03)}, anchor=south east, row sep=1pt},
      legend cell align=left,
      mark size=1.7pt, clip=false,
    ]
    \addplot[cRGB, line width=1.0pt, mark=*, mark options={fill=cRGB, draw=none}]
      coordinates {(1.7,40.9) (2.5,48.6) (4.7,53.1) (6.2,55.0) (11.8,57.5)};
    \addlegendentry{YOLO26 (2026)}
    \addplot[cTime, line width=1.0pt, mark=square*, mark options={fill=cTime, draw=none}]
      coordinates {(3.66,49.8) (5.91,53.7) (8.07,55.4) (12.90,57.0)};
    \addlegendentry{RT-DETRv4 (2025)}
    \addplot[cIR, line width=1.0pt, mark=triangle*, mark options={fill=cIR, draw=none}]
      coordinates {(1.97,41.6) (2.98,48.0) (8.63,53.4) (14.67,54.8)};
    \addlegendentry{YOLOv13 (2025)}
    \addplot[cNeutral, line width=1.0pt, mark=diamond*, mark options={fill=cNeutral, draw=none}]
      coordinates {(1.5,39.5) (2.5,47.0) (4.7,51.5) (6.2,53.4) (11.3,54.7)};
    \addlegendentry{YOLO11 (2024)}
    \addplot[only marks, cFuse, mark=pentagon*, mark size=2.3pt, mark options={fill=cFuse, draw=none}]
      coordinates {(2.3,48.4) (17.2,60.1)};
    \addlegendentry{RF-DETR N and 2XL (2025)}
    \addplot[only marks, cAlert, mark=*, mark size=1.9pt, mark options={fill=white, draw=cAlert, line width=0.8pt}]
      coordinates {(1.60,40.4) (10.38,55.4)};
    \addlegendentry{YOLOv12 N and X (2025)}

    \node[font=\scriptsize\itshape, text=muted, anchor=west] at (axis cs:12.1,57.9) {YOLO26-x};
    \node[font=\scriptsize\itshape, text=muted, anchor=east] at (axis cs:16.9,60.6) {RF-DETR-2XL};
  \end{axis}
\end{tikzpicture}

\vspace{1mm}
\caption[Accuracy against latency for current real-time detectors]{Accuracy
against latency for current real-time detectors, using the figures each work
reports~\citep{yolo26_2026,rtdetrv4_2025,lei2025yolov13,yolo11_2024,robinson2025rfdetr,tian2025yolov12}.
Lines join model sizes whose individual latencies were confirmed. For
RF-DETR and YOLOv12, only the smallest and largest models are plotted.
Protocols differ, notably in whether NMS is timed, so small horizontal
differences should not be read as significant.}
\label{fig:frontier}
\end{figure}


For this project, YOLO26 is the natural base. It is the newest YOLO and is
maintained in the Ultralytics library~\citep{yolo26_2026}. Its label
assignment already targets small objects. Its modules are defined in a
configuration file, which makes new blocks easy to register. The DETR
frontier is stronger at equal latency, but it is heavier to train on a T4
GPU. It is better treated as a reference point than as a base.

\section{Vision state-space models and Mamba detectors}
\label{sec:sota-ssm}

Vision state-space models are now competitive with Swin and ConvNeXt as
general backbones. Examples include VMamba, Spatial-Mamba, DAMamba,
A2Mamba and PRISMamba~\citep{liu2024vmamba,xiao2025spatialmamba,damamba2025,a2mamba2025,prismamba2026}.
Inside real-time YOLO detectors the evidence is weaker. Mamba YOLO, the
AAAI~2025 baseline that most later work builds on, reaches 52.1 COCO AP with
its large model. YOLO11-l reaches 53.4 at similar
computation~\citep{wang2025mambayolo,yolo11_2024}. The closest work to the
project title is MambaPSA (July 2026). It replaces YOLO26's attention block
with a Mamba block and changes AP by $-0.1$ to $+0.9$ on
VOC~\citep{mambapsa2026}. Research effort in 2026 has moved to the scan
order itself (Table~\ref{tab:sota-ssm}).

\begin{table}[htbp]
\centering
\caption[Vision state-space models and Mamba detectors, newest first]{Vision
state-space models and Mamba detectors, newest first.}
\label{tab:sota-ssm}
\small
\begin{tabularx}{\textwidth}{@{}P{1.35cm}P{3.0cm}P{2.1cm}L@{}}
\hd{Date} & \hd{Work} & \hd{Venue} & \hd{Contribution and result} \\ \gap
09/2026 & YOLO12-MambaScan~\citep{yolo12mambascan2026} & arXiv & YOLO12 with Mamba context; VisDrone at 960\,px, 38.6 AP \\
09/2026 & ScopeMamba-YOLO~\citep{scopemamba2026} & arXiv & off-path selective scan; VisDrone 50.8 AP\textsubscript{50} with 3.6M params \\
07/2026 & MambaPSA~\citep{mambapsa2026} & arXiv & Mamba in place of YOLO26's C2PSA; 12\% fewer FLOPs, $-0.1$ to $+0.9$ AP \\
07/2026 & VNCT~\citep{vnct2026} & arXiv & non-causal trapezoidal SSM; argues directional scans are unnecessary \\
06/2026 & AKCMamba-YOLO~\citep{akcmamba2026} & CVPR 2026 & channel-interaction Mamba modules in YOLO \\
05/2026 & C-GSPN~\citep{cgspn2026} & arXiv & notes that SSMs weaken 2D structure; $4\times$ faster at 2K resolution \\
03/2026 & SF-Mamba~\citep{sfmamba2026} & arXiv & reports Mamba is slow at short token lengths; batch folding as remedy \\
02/2026 & PRISMamba~\citep{prismamba2026} & ICML 2026 & partial ring scan; 84.5\% ImageNet top-1 at 3.9 GFLOPs \\
07/2025 & A2Mamba~\citep{a2mamba2025} & arXiv & $+1.2$ box AP over MambaVision-B with 40\% fewer parameters \\
06/2025 & MambaNeXt-YOLO~\citep{mambanextyolo2025} & arXiv & CNN--Mamba block; runs on Jetson Xavier and Orin NX \\
02/2025 & DAMamba~\citep{damamba2025} & NeurIPS 2025 & dynamic adaptive scan; 48.5 to 50.6 box AP, Mask R-CNN $1\times$ \\
10/2024 & Spatial-Mamba~\citep{xiao2025spatialmamba} & ICLR 2025 & structure-aware state fusion; 47.6 to 50.4 box AP \\
07/2024 & MambaVision~\citep{hatamizadeh2025mambavision} & CVPR 2025 & source paper; Section~\ref{sec:src-mv} \\
06/2024 & Mamba YOLO~\citep{wang2025mambayolo} & AAAI 2025 & Ultralytics-based baseline; COCO 44.5 to 52.1 AP \\
05/2024 & MambaOut~\citep{yu2025mambaout} & CVPR 2025 & SSM unnecessary for classification; may help detection \\
05/2024 & MLLA~\citep{han2024mlla} & NeurIPS 2024 & Mamba as gated linear attention; forget gate matters most \\
01/2024 & VMamba~\citep{liu2024vmamba} & NeurIPS 2024 & four-direction SS2D scan \\
01/2024 & Vim~\citep{zhu2024vim} & ICML 2024 & bidirectional scan \\
\end{tabularx}
\end{table}

\section{Small and tiny objects}
\label{sec:sota-small}

Small-object results are hard to compare because three choices vary between
papers: the input resolution (640, 960 or about 1\,333 pixels), the split
(validation or test-dev) and the metric (AP\textsubscript{50} or COCO AP).
Published VisDrone figures range from 26 to 39 AP, and up to 60
AP\textsubscript{50} at high resolution, largely for these reasons.
Table~\ref{tab:sota-small} reports each result with its protocol where the
paper states it.

\begin{table}[htbp]
\centering
\caption[Small and tiny object detection, newest first]{Small and tiny object
detection, newest first.}
\label{tab:sota-small}
\small
\begin{tabularx}{\textwidth}{@{}P{1.35cm}P{2.75cm}P{2.1cm}L@{}}
\hd{Date} & \hd{Method} & \hd{Venue} & \hd{Reported result} \\ \gap
06/2026 & SFDNet~\citep{sfdnet2026} & ECCV 2026 & spectrum-aware; SODA-D 34.2 AP, SODA-A 39.2 AP, AI-TOD 31.7 AP \\
06/2026 & DERNet~\citep{dernet2026} & arXiv & frequency domain; beats YOLO11 with one sixth of the parameters \\
03/2026 & O2-DEIM~\citep{o2deim2026} & IEEE TGRS & oriented real-time DETR; DOTA-v1.0 80.15 AP\textsubscript{50} at 119 FPS \\
01/2026 & EFSI-DETR~\citep{efsidetr2026} & arXiv & VisDrone 33.1 AP, 24.8 AP\textsubscript{S}\unv \\
01/2026 & TOLF~\citep{tolf2026} & arXiv & normalising flows; $+1.2$ AP over DINO on AI-TOD \\
05/2025 & Dome-DETR~\citep{domedetr2025} & ACM MM 2025 & AI-TOD-V2 34.7 AP; VisDrone 39.0 AP at about 376 GFLOPs \\
01/2025 & Strip R-CNN~\citep{yuan2026stripr} & AAAI 2026 & DOTA-v1.0 82.75 mAP with 30M params \\
01/2025 & UAV-DETR~\citep{uavdetr2025} & arXiv & $+3.1$ AP over RT-DETR on VisDrone \\
04/2024 & DQ-DETR~\citep{huang2024dqdetr} & ECCV 2024 & AI-TOD-V2 30.2 AP \\
2021 & NWD~\citep{wang2021nwd} & ISPRS J. 2022 & Wasserstein matching for tiny boxes \\
\end{tabularx}
\end{table}

\section{Visible--thermal fusion}
\label{sec:sota-rgbt}

Fusion has followed the same arc as general detection. Attention-based
fusion came first: CFT and ICAFusion~\citep{qingyun2021cft,shen2024icafusion}.
State-space fusion followed in 2024 and 2025. Its main members are
Fusion-Mamba, COMO, WaveMamba and MS2Fusion~\citep{dong2025fusionmamba,como2024,wavemamba2025,ms2fusion2025}.
In 2026, methods built on foundation models appeared. RegisterBridgeMM, for
example, uses a frozen DINOv3 encoder~\citep{registerbridge2026}.
Table~\ref{tab:sota-rgbt} shows the four standard benchmarks.

LLVIP is saturated. Every recent method scores 97 to 98.4
AP\textsubscript{50}, and a thermal-only model is close behind. DroneVehicle
results are not comparable across papers, because some use horizontal and
some use oriented boxes. M3FD has no official detection split. KAIST is
scored by miss rate, not AP, and has no clear 2025--2026 leader.

\begin{table}[htbp]
\centering
\caption[Visible--thermal fusion detectors, newest first]{Visible--thermal
fusion detectors, newest first. Each cell gives AP\textsubscript{50} and
mAP\textsubscript{50:95}. A dash means not reported.}
\label{tab:sota-rgbt}
\small
\begin{tabularx}{\textwidth}{@{}P{1.3cm}LR{1.7cm}R{1.7cm}R{1.7cm}R{1.85cm}@{}}
\hd{Date} & \hd{Method} & \hd{LLVIP} & \hd{M3FD} & \hd{FLIR-aligned} & \hd{DroneVehicle} \\ \gap
08/2026 & RegisterBridgeMM~\citep{registerbridge2026} & 98.4 / 70.5 & 92.6 / 64.9 & 88.9 / 49.8 & 81.5 / 61.5 \\
08/2026 & CFGPNet~\citep{cfgpnet2026} & 97.8 / 68.9 & 89.9 / 63.4 & 80.7 / 45.0 & -- \\
11/2025 & MambaRefine-YOLO~\citep{mambarefine2025} & -- & -- & -- & 83.2 AP\textsubscript{50} \\
07/2025 & WaveMamba~\citep{wavemamba2025} & 98.3 / 66.0 & 92.1 / 64.4 & 88.4 / 48.1 & 79.8 / 60.5 \\
07/2025 & MS2Fusion~\citep{ms2fusion2025} & reported & reported & reported & -- \\
12/2024 & COMO~\citep{como2024} & 97.2 / 65.3 & -- & -- & 86.1 / 65.5 \\
04/2024 & Fusion-Mamba~\citep{dong2025fusionmamba} & 97.0 / 64.3 & 88.0 / 61.9 & 84.9 / 47.0 & -- \\
04/2024 & CFMW~\citep{li2025cfmw} & 98.8 / 69.8\unv & -- & -- & -- \\
08/2023 & ICAFusion~\citep{shen2024icafusion} & -- & 81.5 / 54.5 & -- & -- \\
11/2021 & CFT~\citep{qingyun2021cft} & 97.5 / 63.6 & -- & -- & 84.3 / 61.9 \\
\end{tabularx}
\end{table}

\section{Video, tracking and foundation models}
\label{sec:sota-video}

The standard tracking benchmarks MOT17 and MOT20 have levelled off at about
64 to 66 HOTA. New work reports mainly on DanceTrack and SportsMOT, which test
association under uniform appearance and fast motion (Table~\ref{tab:sota-mot}).
SambaMOTR is the only state-space tracker among the leaders. It keeps a
synchronised SSM memory per track~\citep{segu2025sambamotr}.

For video object detection on ImageNet VID, no strong 2025--2026 method was
found besides TMambaDet~\citep{tmambadet2026}. TMambaDet uses a bidirectional,
and therefore offline, temporal Mamba.

Foundation models now reach the detection task directly. SAM~3 detects,
segments and tracks objects named by a text prompt~\citep{sam3_2025}. YOLOE
brings open-vocabulary prompts to the YOLO architecture~\citep{wang2025yoloe}.
DINOv3 has become the backbone or teacher of choice for detection
transformers~\citep{simeoni2025dinov3}.

\begin{table}[htbp]
\centering
\caption[Multi-object tracking, newest first]{Multi-object tracking, newest
first. HOTA on DanceTrack and SportsMOT test sets as reported.}
\label{tab:sota-mot}
\small
\begin{tabularx}{\textwidth}{@{}P{1.35cm}P{2.9cm}P{2.1cm}R{1.85cm}R{1.85cm}L@{}}
\hd{Date} & \hd{Method} & \hd{Venue} & \hd{DanceTrack} & \hd{SportsMOT} & \hd{Note} \\ \gap
07/2026 & SpikingMOT~\citep{spikingmot2026} & arXiv & 56.5 & 74.9 & spiking network \\
06/2026 & SAM~3 mask propagation~\citep{smp2026} & arXiv & -- & 87.2\unv & offline, training-free \\
09/2025 & MATR~\citep{matr2025} & arXiv & 71.3 & 72.7 & with extra data \\
03/2024 & MOTIP~\citep{gao2025motip} & CVPR 2025 & 67.5 & 71.9 & tracking as ID prediction \\
10/2024 & SambaMOTR~\citep{segu2025sambamotr} & ICLR 2025 & 67.2 & 69.8 & state-space track memory \\
11/2022 & MOTRv2~\citep{zhang2023motrv2} & CVPR 2023 & 69.9\unv & -- & YOLOX proposals \\
03/2022 & OC-SORT~\citep{cao2023ocsort} & CVPR 2023 & 55.1\unv & 68.1 & motion-based \\
10/2021 & ByteTrack~\citep{zhang2022bytetrack} & ECCV 2022 & 47.7\unv & -- & MOT17 HOTA 63.1 \\
\end{tabularx}
\end{table}

\section{Caveats that govern fair comparison}
\label{sec:caveats}

Seven problems recur across the literature. The experimental protocol in
Chapter~\ref{ch:protocol} is designed around them.

\begin{enumerate}
  \item \emph{Mamba is not automatically fast.} Its scan beats
        FlashAttention-2 only beyond about 2\,000 tokens~\citep{gu2023mamba,dao2023flash2}.
        YOLO's $P_4$ and $P_5$ maps are shorter than that.
  \item \emph{Installation.} The official \texttt{mamba\_ssm} kernels
        support Linux only~\citep{mambassm_repo}. On Windows, the choices are
        WSL2, a cloud notebook, or a pure-PyTorch scan that is about twice
        as slow~\citep{mambapy_repo}.
  \item \emph{Scan order.} A one-dimensional causal scan breaks
        two-dimensional adjacency. At least five papers in 2026 propose
        remedies, so a fixed order needs justification.
  \item \emph{Baseline drift.} Published numbers for the same baseline differ
        between papers. YOLOv12-N, for example, appears as 40.6, 40.4 and
        40.1 AP. Ultralytics itself recommends YOLO11 or YOLO26 over YOLO12
        for stability~\citep{yolo26_2026}.
  \item \emph{Latency protocol.} Many papers omit NMS from latency. GPU
        thermal throttling can shift measurements by about 25\%. Naive FP16
        export can collapse accuracy. RF-DETR's protocol addresses all three:
        it includes NMS, pauses between passes and measures accuracy on the
        exported engine~\citep{robinson2025rfdetr}.
  \item \emph{Saturated benchmarks.} COCO contains label errors that
        COCO-ReM corrects~\citep{singh2024cocorem}. LLVIP is effectively
        solved. RF100-VL shows that COCO gains often fail to transfer to new
        domains~\citep{robicheaux2025rf100vl}.
  \item \emph{Metric choice.} Small-object papers often report
        AP\textsubscript{50} at high resolution. AP, AP\textsubscript{50} and
        size-specific AP should all be reported, at a fixed resolution and
        against baselines trained with the same schedule.
\end{enumerate}

\section{Implications for this project}

The survey settles four choices.
\begin{itemize}
  \item The base detector should be YOLO26. A DETR reference point should be
        kept for honesty.
  \item State-space blocks should be placed where they are efficient and
        where small objects live, at $P_3$ and below, or along the time axis,
        where sequences are long.
  \item Fusion results should be reported on benchmarks that are not
        saturated, with mAP\textsubscript{50:95}.
  \item The decisive comparison is against temporal and fusion baselines
        from 2025 and 2026, not against YOLOv8.
\end{itemize}
